\chapter{Bellman's original resource and search constructions}
\label{sec:p17-bellman-original-constructions}

This chapter supplies the remaining six presentations in Bellman's
1952 note: the shrinking-resource equation, probability-vector
stopping, the two-resource recurrence, convex splitting, box search,
and the scalar threshold equation
\cite[Theorems~2--7, pp.~717--719]{Bellman1952DynamicProgramming}.
The first bounded-transformation theorem is
\cref{thm:p17-bellman-bounded-transformations}.
The continuous state domains below remain continuous. Each
construction states the solution class and the boundary condition
that make its conclusion valid.

\section{Shrinking resources: the exact continuity condition}

\begin{hbdefinition}{def:p17-bellman-shrinking-model}{Allocations, shrinking state and finite-horizon values}
Fix \(x_0>0\), an integer \(m\geq1\), and \(0<c<1\).
Let
\[
 D=\{\mathbf z\in[0,\infty)^m:\ \textstyle\sum_i z_i\leq x_0\},
 \qquad D(x)=\{\mathbf z\in D:\textstyle\sum_i z_i=x\}.
\]
The reward \(a:D\to[0,\infty)\) and successor
\(b:D\to[0,\infty)\) are continuous, with
\[
                a(\mathbf0)=0,\qquad
                0\leq b(\mathbf z)\leq c\sum_i z_i.
\]
At resource level \(x\), choose \(\mathbf z\in D(x)\),
receive \(a(\mathbf z)\), and continue at \(b(\mathbf z)\).
Put
\[
\begin{aligned}
 (Tf)(x)&=\max_{\mathbf z\in D(x)}
                         [a(\mathbf z)+f(b(\mathbf z))],\\
 F_0(x)&=0,\qquad F_{n+1}=TF_n,\qquad
 V(x)=\sup_{n\geq0}F_n(x)\in[0,\infty],\\
 h(x)&=\max_{\mathbf z\in D(x)}a(\mathbf z),\qquad
 H(t)=\max_{\substack{\mathbf z\in D\\\sum_i z_i\leq t}}a(\mathbf z)
                                           \quad(0\leq t\leq x_0).
\end{aligned}
\]
The maximum defining \(T\) is used for continuous \(f\).
\(F_n\) is the optimal total reward for \(n\) decisions
with terminal reward zero. Infinite total reward means
the nonnegative series along a trajectory; \(V\) is its
supremum over trajectories. There is no time discount.
\end{hbdefinition}

\begin{hbtheorem}{thm:p17-bellman-shrinking}{proved=annex:p17.bellman.shrinking}{The continuous solution and a constructive sufficient condition}
Every \(F_n\) is continuous, \(F_n(0)=0\), and \(F_n\)
increases with \(n\). The equation \(F=TF\) has a continuous
solution on \([0,x_0]\) with \(F(0)=0\) if and only if
\[
             \lim_{\delta\downarrow0}
                    \sup_{0\leq x\leq\min(\delta,x_0)}V(x)=0.
\]
When this condition holds, the solution is unique and equals \(V\).
Writing \(\eta(t)=\sup_{0\leq x\leq t}V(x)\), one has
\[
                 0\leq V(x)-F_n(x)\leq\eta(c^n x_0).
\]
In particular the sufficient envelope condition
\[
                           \sum_{j=0}^{\infty}H(c^j x_0)<\infty
\]
ensures existence and gives the uniform tail bound
\[
          0\leq V(x)-F_n(x)
                          \leq\sum_{j=n}^{\infty}H(c^j x_0).
\]
The original geometric values \(h(c^j x_0)\) alone do not
give this conclusion without further control between those levels.

For linear data \(a(\mathbf z)=\sum_i r_i z_i\),
\(b(\mathbf z)=\sum_i d_i z_i\), with \(r_i\geq0\) and
\(0\leq d_i\leq c\), the exact solution is
\[
                     V(x)=Kx,\qquad
                     K=\max_i\frac{r_i}{1-d_i}.
\]
An index attaining this maximum can receive the entire
resource at every stage. If \(x_0=0\), the domain consists
only of zero and the boundary condition gives the sole value zero.
\end{hbtheorem}

\begin{hbtheorem}{thm:p17-bellman-shrinking-counterexample}{proved=annex:p17.bellman.shrinking.counterexample}{A continuous counterexample to geometric-grid summability}
Take \(m=1,x_0=1,c=1/2\), \(b(x)=x/4\), and
\[
 a(x)=\sum_{j=0}^{\infty}\frac1{j+1}
       \left(1-8\,4^j\left|x-\frac34\,4^{-j}\right|\right)_+,
             \qquad (t)_+=\max(t,0).
\]
This is a nonnegative continuous reward on \([0,1]\)
with \(a(0)=0\), and
\(\sum_{j\geq0}h(2^{-j})=0\).
Nevertheless the total reward from \(x=3/4\) is infinite.
There is no continuous solution with boundary value zero.
\end{hbtheorem}

\section{The probability-vector stopping equation}

\begin{hbdefinition}{def:p17-bellman-probability-stopping}{Absorption, transformations and a probe}
Let \(d\geq1\), \(M\geq1\), and
\[
 \Delta_d=\{(p_0,\ldots,p_d):p_i\geq0,\ \textstyle\sum_i p_i=1\}.
\]
Write \(e_i\) for the coordinate probability vector
(the source's \(x_i\)); \(e_0\) is absorbing and has cost zero.
At any other vector one pays one unit and chooses either:
a probe, whose successor is \(e_i\) with probability \(p_i\);
or one of the maps \(S_l:\Delta_d\to\Delta_d\), \(1\leq l\leq M\).
Assume constants \(0<c_l<1\) satisfy
\[
              \sum_{i=1}^d(S_lp)_i
                            \leq c_l\sum_{i=1}^d p_i.
\]
No continuity of \(S_l\) is assumed.
The value is the infimum expected number of steps to absorption,
over all policies on the successive histories; failure ever to
absorb has infinite step cost.
For a fixed initial vector the histories form a finitely branching
tree, so all probabilities used here are defined by its finite
products. No global Borel regularity of the maps is being asserted.
\end{hbdefinition}

\begin{hbtheorem}{thm:p17-bellman-probability-stopping}{proved=annex:p17.bellman.probability.stopping}{Existence, uniqueness, policy and approximation on the entire simplex}
There is a unique bounded real solution of
\[
 \begin{split}
 F(e_0)&=0,\\
 F(p)&=\min\left\{1+\sum_{i=0}^d p_iF(e_i),\
                          1+F(S_1p),\ldots,1+F(S_Mp)\right\}
                                                     \quad(p\ne e_0).
 \end{split}
\]
It is the stated minimum expected absorption time, is at least
one off \(e_0\), and is attained by choosing a minimizing
action at each vector.
For any selected \(l\), let \(K=2/(1-c_l)\). Then \(F\leq K\).
Starting with \(F_0=0\), applying the displayed operator
with value zero fixed at \(e_0\) gives
\[
                0\leq F(p)-F_n(p)\leq K^2/n\qquad(n\geq1).
\]
Thus successive approximation is justified uniformly, without
assuming a contraction with a time discount or continuity at \(e_0\).
If the strict mass-reduction premise is dropped, \(d=M=1\)
with the identity map has no finite solution at \(e_1\).
\end{hbtheorem}

\section{Two resources and the priority comparison}

\begin{hbdefinition}{def:p17-bellman-two-resources}{The original recurrence and its value class}
Fix \(0<p_1,p_2,s_1,s_2<1\) and \(r_1,r_2>0\).
At \((x,y)\in[0,\infty)^2\), action 1 gives
the expression \(p_1[r_1x+F(s_1x,y)]\), and action 2
gives \(p_2[r_2y+F(x,s_2y)]\).
For an infinite action sequence \(w_0,w_1,\ldots\in\{1,2\}\),
its return is
\[
 \sum_{j\geq0}\left(\prod_{h=0}^j p_{w_h}\right)
 \begin{cases}
      r_1s_1^{N_1(j)}x,&w_j=1,\\
      r_2s_2^{N_2(j)}y,&w_j=2,
 \end{cases}
\]
where \(N_i(j)\) counts the occurrences of \(i\) before position \(j\).
Define \(V(x,y)\) as the supremum of these returns.
All such series converge on each bounded rectangle.
The solution class for the recurrence is the functions bounded
on every rectangle \([0,X]\times[0,Y]\).
\end{hbdefinition}

\begin{hbtheorem}{thm:p17-bellman-two-resources}{proved=annex:p17.bellman.two.resources}{The complete priority rule, ties and axes}
The value \(V\) is the unique solution in the stated class of
\[
           V(x,y)=\max\{p_1[r_1x+V(s_1x,y)],
                             p_2[r_2y+V(x,s_2y)]\}.
\]
It is continuous, convex and positively homogeneous.
Define the two current indices
\[
                         I_1=\frac{p_1r_1x}{1-p_1},\qquad
                         I_2=\frac{p_2r_2y}{1-p_2}.
\]
If \(I_1\geq I_2\), action 1 attains the maximum;
if \(I_2\geq I_1\), action 2 does. Thus both are optimal
at equality, and choosing a larger current index at every
stage attains \(V\).
The boundary values are
\[
     V(0,0)=0,\qquad
     V(x,0)=\frac{p_1r_1x}{1-p_1s_1},\qquad
     V(0,y)=\frac{p_2r_2y}{1-p_2s_2}.
\]
For any fixed rectangle, finite-horizon iteration from zero
has a geometric tail bounded by
\[
 \frac{\bar p^{\,n}}{1-\bar p}(p_1r_1X+p_2r_2Y),
                         \qquad \bar p=\max(p_1,p_2)<1.
\]
\end{hbtheorem}

\begin{hbtheorem}{thm:p17-bellman-commensurate}{proved=annex:p17.bellman.commensurate}{The actual piecewise-linear construction}
Suppose \(s_1^m=s_2^n\) for positive integers \(m,n\).
Choose relatively prime positive integers \(u,v\) and
\(0<q<1\) so that \(s_1=q^u,s_2=q^v\).
For \(x,y>0\), on each open cone
\[
                       q^{k+1}<I_1/I_2<q^k,\qquad k\in\mathbb Z,
\]
the optimal action sequence is determined by the integer recursion
\[
       k\longmapsto
       \begin{cases}k+u,&k<0\quad(\hbox{action 1}),\\
                     k-v,&k\geq0\quad(\hbox{action 2}).\end{cases}
\]
It enters \(\{-v,\ldots,u-1\}\) after finitely many steps.
There it is periodic with period \(u+v\), containing
\(v\) actions of type 1 and \(u\) of type 2.

For a finite word \(w\), let \(D_w=\prod_jp_{w_j}\),
let \(n_i(w)\) be its number of type-\(i\) actions, and define
\[
 A_w=r_1\sum_{j:w_j=1}\left(\prod_{h=0}^j p_{w_h}\right)s_1^{N_1(j)},
 \qquad
 B_w=r_2\sum_{j:w_j=2}\left(\prod_{h=0}^j p_{w_h}\right)s_2^{N_2(j)}.
\]
Empty words have \(D_w=1,A_w=B_w=0,n_i(w)=0\).
If \(w\) is the entrance prefix and \(z\) the subsequent
cycle, then the exact value on that cone is
\[
 V(x,y)=A_wx+B_wy+
 D_w\frac{A_zs_1^{n_1(w)}x+B_zs_2^{n_2(w)}y}
                     {1-p_1^v p_2^u q^{uv}}.
\]
This gives a linear function on each of countably many cones,
locally finitely many away from the axes. Adjacent formulas
agree at their common boundary by continuity; the axes have
the preceding explicit values.

For \(p_1=p_2=s_1=s_2=1/2\), \(r_1=r_2=1\),
\[
             V(x,y)=\frac{4x+2y}{7}\quad(1\leq x/y\leq2),
 \qquad
             V(x,y)=\frac{2x+4y}{7}\quad(1/2\leq x/y\leq1).
\]
In particular \(V(3/2,1)=8/7\).
\end{hbtheorem}

\section{Convex splitting and its concave boundary}

\begin{hbtheorem}{thm:p17-bellman-convex-split}{proved=annex:p17.bellman.convex.split}{The continuous split problem reduces to endpoints}
Fix \(x_0>0\), \(0<a,b<1\), and continuous,
nondecreasing convex functions \(g,h:[0,x_0]\to\mathbb R\)
with \(g(0)=h(0)=0\). This includes the source's
twice-differentiable case \(g'',h''\geq0\); continuity
of the second derivatives is not required.
The unique continuous value with \(F(0)=0\) for
\[
 F(x)=\max_{0\leq y\leq x}
          \{g(y)+h(x-y)+F(ay+b(x-y))\}
\]
is convex and satisfies exactly
\[
                   F(x)=\max\{g(x)+F(ax),\ h(x)+F(bx)\}.
\]
Thus an optimal split assigns the entire resource to one
of the two activities. Put
\[
 c=\max(a,b),\qquad K=\max(g(x_0),h(x_0))/x_0.
\]
The \(n\)-step endpoint recursion from zero differs from
\(F(x)\) by at most \(Kxc^n/(1-c)\).
Equivalently, \(F(x)\) is the supremum of the convergent
reward series over all binary endpoint sequences, with that
explicit uniform tail bound.

For \(g(x)=Ax^\lambda,h(x)=Bx^\lambda\), \(A,B\geq0\)
and \(\lambda\geq1\), the same result holds directly under
convexity and gives
\[
        F(x)=
          \max\!\left\{\frac A{1-a^\lambda},
                       \frac B{1-b^\lambda}\right\}x^\lambda .
\]
For example \(g=x^2,h=2x^2,a=1/2,b=1/4\) gives
\(F(x)=32x^2/15\), attained by the second activity.
\end{hbtheorem}

\begin{hbtheorem}{thm:p17-bellman-concave-split}{proved=annex:p17.bellman.concave.split}{A complete interior-split counterexample}
For \(x_0=1\), \(a=b=1/2\), and
\(g(x)=h(x)=x-x^2/2\), the functions are increasing
and strictly concave. The unique continuous value is
\[
                          F(x)=2x-x^2/3.
\]
At every \(x>0\) the optimal split is \(y=x/2\).
The endpoint-only problem instead has value
\(2x-2x^2/3\), strictly smaller for every \(x>0\).
\end{hbtheorem}

\section{Box search: retrieval and location are different tasks}

\begin{hbdefinition}{def:p17-bellman-box-search}{Boxes, completed inspections and attempt costs}
Exactly one of \(d\geq1\) boxes contains the ball.
Its probabilities are \(p_i\geq0,\sum_i p_i=1\).
Here indices range from 1 to \(d\),
\(\operatorname{supp}p=\{i:p_i>0\}\), and \(e_i\)
denotes the unit prior on box \(i\).
An attempt on box \(i\) costs \(t_i>0\).
Conditional on the past and on the ball's location, an
attempt fails to inspect its contents with probability
\(q_i\in[0,1)\); otherwise it reveals whether the ball is
there and retrieves it if present.
This conditional failure law is the operational premise
for the source's completed-inspection cost
\[
                              c_i=\frac{t_i}{1-q_i}.
\]
A failed attempt supplies no location information.
After a successful empty inspection of \(i\), the
remaining probabilities are \(p_j/(1-p_i)\), \(j\ne i\).
Denote that prior by \(p^{(-i)}\), assigning coordinate
zero to the inspected box.
If \(p_i=1\), there is no remaining-state distribution:
the continuation term is zero, without evaluating that quotient.
Zero-probability boxes can be ignored.
\end{hbdefinition}

\begin{hbtheorem}{thm:p17-bellman-box-retrieval}{proved=annex:p17.bellman.box.retrieval}{The full ordering and the original recursion}
Let \(\sigma\) order the positive-probability boxes so that
\[
            \frac{p_{\sigma(1)}}{c_{\sigma(1)}}\geq
            \frac{p_{\sigma(2)}}{c_{\sigma(2)}}\geq\cdots .
\]
The minimum expected retrieval cost, among all randomized
history-dependent attempt policies, is
\[
 F(p)=\sum_{j=1}^{|\operatorname{supp}p|}
             c_{\sigma(j)}
                    \left(1-\sum_{h<j}p_{\sigma(h)}\right).
\]
It is attained by completing inspections in this order,
repeating the current box after a failed attempt.
Ties may be ordered arbitrarily.
For every probability vector it satisfies the original equation
\[
 F(p)=\min_{i:p_i>0}
           \{c_i+(1-p_i)F(p^{(-i)})\},\qquad F(e_i)=c_i,
\]
with the zero continuation convention at \(p_i=1\).
The source's index is exactly \(p_i(1-q_i)/t_i\).
The formula is valid on the entire probability simplex
and is linear on each region of a fixed ordering.

If a positive-probability box has \(q_i=1\), retrieval
cannot occur almost surely and its expected completion
time is infinite. A box with \(p_i=0\) remains irrelevant.
The equation as an equation of completed costs does not
need an attempt model; using \(t_i/(1-q_i)\) as that cost
requires the stated conditional failure law.
If zero attempt costs are allowed, the same formula holds
among policies required to complete retrieval almost surely:
put every positive-probability zero-cost box first and use
the equivalent cross-product comparisons \(p_ic_j\geq p_jc_i\).
The printed ratio index is used only when its denominator is positive.
\end{hbtheorem}

\begin{hbtheorem}{thm:p17-bellman-box-location}{proved=annex:p17.bellman.box.location}{The source's location-only qualification, completed}
If the task is only to identify the occupied box, the last
remaining positive-probability box need not be inspected.
For each proposed uninspected last box \(l\), order all
other positive-probability boxes by decreasing \(p_i/c_i\)
and compute
\[
 L_l(p)=\sum_{j=1}^{m-1}c_{\sigma(j)}
                 \left(1-\sum_{h<j}p_{\sigma(h)}\right),
                     \qquad m=|\operatorname{supp}p|,\quad \sigma(m)=l.
\]
Then \(L(p)=\min_l L_l(p)\) is the minimum expected
location cost. For any fixed highest-index box there
is an optimal policy which inspects it first or leaves
it uninspected. If \(m=1\), \(L=0\), whereas retrieval
still costs \(c_i\).
If inspection probabilities \(q_i=1\) are allowed, location
has finite cost exactly when at most one positive-probability
box is uninspectable. When there is one, it must be the
uninspected last box and the same sorted-prefix formula applies
to the other boxes.

For \(p=(1/2,3/10,1/5)\), \(q_i=0\) and
\(t=(2,1,4)\), the retrieval order is \(2,1,3\),
with \(F=16/5\); location omits box 3 and costs \(12/5\).
For \(p=(4/5,1/5)\), \(t=(1,3/10)\), \(q_i=0\),
retrieval starts with the highest-index box 1 and costs
\(53/50\); location leaves that box uninspected and costs \(3/10\).
\end{hbtheorem}

\section{The scalar threshold equation}

\begin{hbtheorem}{thm:p17-bellman-threshold}{proved=annex:p17.bellman.threshold}{The exact threshold, minimizing integer and all ties}
Fix \(0<a<1\). The unique bounded solution with \(f(0)=0\) of
\[
                   f(x)=1+\min\{x f(1),\,f(ax)\},
                                    \qquad 0<x\leq1,
\]
is obtained as follows:
\[
 C=\min_{k\geq1}\frac{k+1}{1-a^k},\qquad
 \tau=\frac1{C(1-a)}.
\]
The integer minimum is attained; it suffices to examine
\(1\leq k\leq\lfloor2/(1-a)-1\rfloor\).
Every minimizing \(k\) gives the same threshold
\[
                 \tau=\frac{1-a^k}{(k+1)(1-a)},\qquad 1/2\leq\tau<1.
\]
For \(x>0\), let \(n(x)\) be the least nonnegative integer
with \(a^{n(x)}x\leq\tau\). Then
\[
                f(x)=n(x)+1+a^{n(x)}x C,\qquad f(1)=C.
\]
In particular
\[
 f(x)=1+xC\quad(0<x\leq\tau),\qquad
 f(x)=1+f(ax)\quad(\tau\leq x\leq1).
\]
Both actions are optimal at the threshold. A tie
\(a^n x=\tau\) allows probing after either \(n\) or \(n+1\)
transformations, with the same cost.
The value has limit one as \(x\downarrow0\), while its
absorbing boundary value is \(f(0)=0\).
At the excluded parameter endpoint \(a=0\), the equation
instead has \(f(x)=1\) for \(x>0\); at \(a=1\) it has
no finite solution.

For \(a=1/2\), \(C=4\), both \(k=1,2\) minimize,
\(\tau=1/2\), and
\[
 f(x)=\begin{cases}
 0,&x=0,\\
 1+4x,&0<x\leq1/2,\\
 2+2x,&1/2\leq x\leq1.
 \end{cases}
\]
\end{hbtheorem}
